\documentclass[11pt]{article}
\usepackage[T1]{fontenc}
\usepackage[margin=1in]{geometry}
\usepackage{amsmath,amssymb,amsthm,mathtools}
\usepackage{graphicx}
\usepackage{microtype}
\usepackage{xcolor}
\usepackage[colorlinks=true,linkcolor=blue!55!black,citecolor=blue!55!black,urlcolor=blue!65!black]{hyperref}

\newtheorem{theorem}{Theorem}
\newtheorem{corollary}[theorem]{Corollary}
\newcommand{\SG}{\mathcal{SG}}
\newcommand{\E}{\mathcal E}
\newcommand{\F}{\mathcal F}
\newcommand{\R}{\mathsf R}
\emergencystretch=3em

\title{Finite-energy extension from every level-$l$ half Sierpi\'nski gasket}
\author{Partial-result packet for arXiv:1702.02419}
\date{27 August 2026}

\begin{document}
\maketitle

\begin{center}
\fbox{\begin{minipage}{0.91\textwidth}
\textbf{Status: candidate partial result, likely valid.}
For every $l\geq2$, reflection gives a bounded linear extension of every
finite-energy function on the vertical half domain to the whole $\SG_l$, with
an exact energy identity.  The source's upper/lower-domain and higher-Sobolev
extension problems remain open.
\end{minipage}}
\end{center}

\section{Source problem}

Cao and Qiu study half, upper, and lower domains in the level-$l$ Sierpi\'nski
gaskets and ask on PDF page 6 how to extend a finite-energy function from such
a domain to the whole gasket \cite[p.~6]{CQ}.  Their exact sentence is:
\begin{quote}
``One could also consider how to extend a function of finite energy defined on
$\Omega$ to a function of finite energy on the whole $\SG_l$, and analogous
problems for other Sobolev spaces.''
\end{quote}
The source passage is reproduced below.

\begin{center}
\includegraphics[width=\textwidth]{figures/open_problem_crop.png}
\end{center}

The theorem below resolves the finite-energy question for the full vertical
half-domain family.  It is stronger than a harmonic-extension statement: the
input function is arbitrary in the local resistance-form domain.

\section{Definitions and statement}

Let $\rho$ be reflection in the vertical symmetry axis of $\SG_l$.  Let
$h_a$ be the antisymmetric harmonic function used by Cao--Qiu, normalized by
$h_a(q_0)=0$, $h_a(q_1)=1$, and $h_a(q_2)=-1$.  Put
\[
  \Omega^- = \{h_a>0\},\qquad
  X=\{h_a=0\},\qquad
  \Omega^+=\rho(\Omega^-)=\{h_a<0\}.
\]
Thus $X$ is the pointwise fixed set of $\rho$, and
$\SG_l=\Omega^-\sqcup X\sqcup\Omega^+$.

Write $\Gamma_m=(V_m,E_m)$ for the usual level-$m$ graph approximation and
$r_l$ for the resistance renormalization factor.  The local graph energy on
the closed left half is
\[
 \E_{m,-}(u)=r_l^{-m}
   \sum_{e=\{x,y\}\in E_m^-}(u(x)-u(y))^2,
\]
where $E_m^-$ is the left member of every non-fixed reflection orbit of graph
edges.  Equivalently, these are the edge terms in the standard left-half
exhaustion used for $\E_{\Omega^-}$ in the source.  Let
\[
 \F(\overline{\Omega^-})
 =\{u\in C(\overline{\Omega^-}):
       \E_{\Omega^-}(u):=\lim_m\E_{m,-}(u)<\infty\}.
\]

\subsection*{Proof intuition}

At each graph level, reflection partitions the non-fixed edges into pairs.
An even extension gives the same energy to both edges in a pair.  The only
edges fixed as sets are horizontal edges crossing the symmetry axis; reflection
exchanges their endpoints, so the even extension has equal endpoint values and
their energy is zero.  The full graph energy is therefore exactly twice the
left-half energy at every finite level.  Passing to the resistance-form limit
simultaneously proves existence, boundedness, and the sharp identity for this
specific extension.

\begin{theorem}[even reflection extension]\label{thm:main}
For every integer $l\geq2$, define
\[
  (\R u)(x)=
  \begin{cases}
    u(x),&x\in\overline{\Omega^-},\\
    u(\rho x),&x\in\Omega^+.
  \end{cases}
\]
Then
\[
  \R:\F(\overline{\Omega^-})\longrightarrow\F(\SG_l)
\]
is a linear right inverse to restriction, and
\begin{equation}\label{eq:energy}
  \E_{\SG_l}(\R u)=2\E_{\Omega^-}(u).
\end{equation}
It also preserves the uniform norm:
$\|\R u\|_{\infty,\SG_l}=\|u\|_{\infty,\overline{\Omega^-}}$.
Consequently, with the anchored resistance norm
\[
 \|u\|_{*,Y}^2=|u(q_0)|^2+\E_Y(u),
\]
the extension operator has norm at most $\sqrt2$.
\end{theorem}

\section{Proof}

The reflection $\rho$ permutes the contraction maps defining $\SG_l$ and is
therefore an automorphism of every graph $\Gamma_m$.  Hence it permutes
$E_m$.  The triangular lattice has no nontrivial edge contained in the
vertical axis.  Every edge orbit consequently has one of two forms:
\begin{enumerate}
\item $\{e,\rho e\}$, where $e$ is on the left and $\rho e$ is on the right;
\item a singleton $\{e\}$, where $e=\{x,\rho x\}$ crosses the axis and
      reflection exchanges its endpoints.
\end{enumerate}
This classification also follows cell by cell: the three graph-edge directions
are horizontal and at angles $\pm\pi/3$, whereas the fixed axis is vertical.

Because $\rho$ fixes $X$ pointwise, the two formulas defining $\R u$ agree on
$X$.  They therefore define a continuous function on $\SG_l$.  For a paired
orbit, say $e=\{x,y\}$ on the left, the two contributions to the graph energy
of $\R u$ are equal:
\[
 \bigl((\R u)(x)-(\R u)(y)\bigr)^2
 =\bigl((\R u)(\rho x)-(\R u)(\rho y)\bigr)^2
 =(u(x)-u(y))^2.
\]
For a singleton orbit $e=\{x,\rho x\}$, evenness gives
$(\R u)(x)=(\R u)(\rho x)$, so its contribution is zero.  Summing edge
orbits and inserting the common factor $r_l^{-m}$ gives the exact finite-level
identity
\[
  \E_m(\R u)=2\E_{m,-}(u)
  \qquad(m\geq0).
\]
Taking increasing limits proves \eqref{eq:energy}; in particular $\R u$ has
finite global energy.  Linearity, the restriction identity, and the uniform
norm identity are immediate from the definition.  Finally,
\[
 \|\R u\|_{*,\SG_l}^2
 =|u(q_0)|^2+2\E_{\Omega^-}(u)
 \leq2\bigl(|u(q_0)|^2+\E_{\Omega^-}(u)\bigr),
\]
which proves the stated norm bound. \qed

\begin{corollary}[energy-minimizing completion of the other half]
For each $u\in\F(\overline{\Omega^-})$, there is a global finite-energy
extension that agrees with $u$ on $\overline{\Omega^-}$ and is harmonic on
$\Omega^+$.  Its total energy is at most $2\E_{\Omega^-}(u)$.
\end{corollary}

\begin{proof}
The reflected function $\R u$ supplies a finite-energy competitor on
$\Omega^+$ with boundary trace $u|_X$.  Minimizing the right-half energy with
that trace, by the usual Dirichlet principle for the resistance form (the same
principle used in Proposition 1.1 of \cite{CQ}), gives a harmonic completion
whose right-half energy is no larger than that of the reflected competitor,
namely $\E_{\Omega^-}(u)$.  Add the unchanged left-half energy.
\end{proof}

\section{Verification, scope, and novelty audit}

The edge-orbit identity was checked independently by
\nolinkurl{code/verify_reflection_energy.py}.  The script constructs the
standard triangular IFS graphs for $2\leq l\leq6$, through depth three for
$l\leq4$ and depth two otherwise, verifies that reflection is an edge
automorphism and that no nontrivial edge lies in the fixed axis, and tests the
energy identity on seeded random symmetric vertex data.  This is a finite
regression check only; the edge-orbit argument above is the proof.

The result is deliberately scoped to the vertical half domains and the
first-order resistance energy.  It does not settle arbitrary horizontal upper
or lower cuts, nor extension in $\operatorname{dom}_{L^2}(\Delta^k)$.  Even
reflection also need not make a harmonic function globally harmonic across
$X$: a normal-flux matching condition can fail.

The run searched its registry, solution, attempt, and proof-gap indexes for
arXiv:1702.02419 and the core extension terms.  Bounded primary-source searches
for half-gasket finite-energy reflection/extension recovered Li--Strichartz's
$\SG_2$ work on trace spaces and piecewise-biharmonic extensions for Laplacian
domains \cite{LS}, Cao--Qiu's later Sobolev work on p.c.f. sets \cite{CQ20},
and a bottom-line trace theorem \cite{CLST}, but no all-$l$ statement of
Theorem~\ref{thm:main}.  Because the symmetry argument is elementary, the
novelty claim is plausible but especially provisional: the theorem may be
folklore or implicit in resistance-form locality.

\paragraph{Human-review focus.}
Confirm that the paper's local energy $\E_{\Omega^-}$ agrees with the left-edge
exhaustion used here for all $l$, especially at graph edges setwise fixed by
reflection; then check the bounded literature search for an implicit reflection
extension lemma.

\begin{thebibliography}{9}
\bibitem{CQ}
S.~Cao and H.~Qiu,
\emph{Boundary Value Problems for harmonic functions on domains in Sierpi\'nski
gaskets}, arXiv:1702.02419 (2017).

\bibitem{LS}
W.~Li and R.~S.~Strichartz,
\emph{Boundary Value Problems on a Half Sierpi\'nski Gasket},
J. Fractal Geom. 1 (2014), 1--43; arXiv:1304.3522.

\bibitem{CQ20}
S.~Cao and H.~Qiu,
\emph{Sobolev spaces on p.c.f. self-similar sets: critical orders and atomic
decompositions}, arXiv:1904.00342 (2020).

\bibitem{CLST}
S.~Cao, S.~Li, R.~S.~Strichartz, and P.~Talwai,
\emph{A Trace Theorem For Sobolev Spaces On The Sierpi\'nski Gasket},
arXiv:1905.03391 (2019).
\end{thebibliography}

\end{document}
